\documentclass[runningheads]{llncs}

\usepackage[T1]{fontenc}
\usepackage{amsmath,amssymb}
\usepackage{url}

\begin{document}

\title{Estudio comparativo del rendimiento de un árbol de decisión C4.5 propio y del modelo de ensamble Random Forest en problemas de clasificación}
\titlerunning{C4.5 y Random Forest}

\author{Laura Cornaló Bassan\inst{1}\textsuperscript{[00000]} \and
Matías Fernandez\inst{1}\textsuperscript{[24282]} \and
Mariano Insaurralde\inst{1}\textsuperscript{[24684]} \and
Philippe Maurel Garcete\inst{1}\textsuperscript{[25613]} \and
Franco Sanchez\inst{1}\textsuperscript{[27310]}}
\authorrunning{L. Cornaló Bassan, M. Fernandez, M. Insaurralde,
P. Maurel Garcete, F. Sanchez}

\institute{Universidad Tecnológica Nacional, Facultad Regional de Resistencia\\
Ingeniería en Sistemas de Información\\
\email{\{correo1,correo2,correo3,correo4,correo5\}@gmail.com}}

\maketitle

\begin{abstract}
% TODO: redactar resumen.
\keywords{Decision Tree \and C4.5 \and Random Forest}
\end{abstract}

\section{Introducción}
% TODO: redactar introducción.

\section{Fundamentos teóricos}

\subsection{Árboles de decisión y C4.5}
% TODO: redactar la base teórica y citar \cite{liu2011}.

\subsection{Random Forest}
% TODO: redactar la base teórica y citar \cite{breiman2001,scikitRF}.

\section{Breve descripción de los algoritmos desarrollados}
% TODO: completar cuando estén implementados.

\section{Solución implementada}
% TODO: describir la solución y los problemas encontrados.

\section{Simulaciones y resultados}
% TODO: incorporar resultados de experimentos reales.

\section{Conclusiones}
% TODO: redactar conclusiones después de analizar los resultados.

\nocite{liu2011,breiman2001,scikitRF}
\bibliographystyle{splncs04}
\bibliography{refs}

\end{document}
